\documentclass{report}
\usepackage{amsmath,amssymb}
\setlength{\parindent}{0mm}
\setlength{\parskip}{1em}
\begin{document}
\begin{center}
\rule{6in}{1pt} \
{\large Alessio Spantini \\
{\bf Optimal low-rank approximations of goal-oriented Bayesian linear inverse problems}}

888 MASSACHUSETTS Avenue Apt 318 \\ Cambridge \\ 02139 (MA) \\ USA
\\
{\tt spantini@mit.edu}\\
Cui Tiangang\\
Luis Tenorio\\
Youssef Marzouk\end{center}

We propose statistically optimal and computationally efficient
dimensionality reduction techniques for finite-dimensional goal-oriented
Bayesian Gaussian linear inverse problems where the QoI is a linear
function of the inversion parameters. In particular, we study
approximations of posterior covariance of the QoI as a low-rank negative
update of the prior covariance of the QoI and prove optimality of the
update with respect to the natural geodesic distance on the manifold of
symmetric and positive definite matrices. If we assume exact knowledge of
the posterior mean, then our optimality results extend to optimality in
distribution with respect to the Kullback-Leibler divergence and the
Hellinger distance between the associated distributions. We also propose
approximations of the posterior mean of the QoI as a low-rank linear
function of the data and prove optimality of the approximation with
respect to the Bayes risk for squared-error loss weighted by the
posterior precision matrix of the QoI. These optimal approximations avoid
the explicit computation of the full posterior distribution of the
parameters and focus entirely on directions that are well informed by the
data and that are relevant to the QoI. These directions are obtained from
the leading generalized eigenvectors of a suitable matrix pencil and stem
from a careful balance between all the ingredients of the goal oriented
inverse problem: prior information, forward model, measurement noise and
ultimate goals. We illustrate the theory using a high-dimensional inverse
problem in heat transfer.


\end{document}
